\section{Experimental Setup}

\subsection{Memory Backends}
We evaluate five memory backends spanning different retrieval paradigms.
\textbf{FAISS}~\citep{johnson2019faiss} uses dense vector similarity search with SiliconFlow's multilingual embedding API and top-$k{=}5$ retrieval.
\textbf{BM25} uses sparse keyword matching via rank-BM25 with top-$k{=}5$ retrieval.
\textbf{Naive RAG} uses sequential passage concatenation (no relevance ranking), top-$k{=}5$.
\textbf{Simple Knowledge Graph} is an in-memory graph with entity/relation extraction, returning connected subgraphs (5--10 facts per query).
\textbf{ChromaDB}~\citep{chromadb2023} is a persistent vector database using its default embedding function (all-MiniLM-L6-v2).

All systems share identical chunking, top-$k$ parameters, generation prompts, and the LLM backbone (DeepSeek-V3 via SiliconFlow API), ensuring differences reflect retrieval design choices.

\subsection{Cross-Benchmark Evaluation}
We additionally evaluate all backends (except ChromaDB) on the LoCoMo benchmark~\citep{locomo2024}, which contains 1,986 MC questions across 10 long conversations (300--700 turns each). This enables cross-benchmark analysis of system behavior under different conditions: BiTempQA's short, dual-timestamp scenarios vs.\ LoCoMo's long, single-timestamp conversations.

\subsection{Statistical Methods}
We employ: (1) Bootstrap 95\% confidence intervals (10,000 iterations); (2) McNemar's test with continuity correction for pairwise system comparisons; (3) Mixed-effects logistic regression with scenario random intercepts to control for question type, difficulty, and temporal requirements; (4) LLM Judge cross-validation (DeepSeek-V3 vs.\ Qwen2.5-72B, Cohen's $\kappa{=}0.455$).
